\begin{solution}{96}
    \begin{subsolution}
        (解法一)

        半球和小球组成的系统在水平方向上没有受到外力作用，系统在水平方向上动量守恒
        \begin{equation}
            - M V _ {M} + m [ (R + r) \dot {\theta} \cos \theta - V _ {M} ] = 0
            \label{eq:1.p96}
        \end{equation}
        设小球转动角速度大小为 $\omega$ ，小球做纯滚动，故有
        \begin{equation}
            r \omega = (R + r) \dot {\theta}
            \label{eq:2.p96}
        \end{equation}
        无耗散力做功，系统的机械能守恒
        \begin{equation}
            m g (R + r) (1 - \cos \theta)
            = \frac {1}{2} M V _ {M} ^ {2} + \frac {1}{2} m \left\{\left[ (R + r) \dot {\theta} \cos \theta - V _ {M} \right] ^ {2} + \left[ (R + r) \dot {\theta} \sin \theta \right] ^ {2} \right\} + \frac {1}{2} I \omega^ {2}
            \label{eq:3.p96}
        \end{equation}
        式中 $I=\frac{2}{5}mr^{2}$ 。联立\eqref{eq:1.p96}\eqref{eq:2.p96}\eqref{eq:3.p96}式得，小球运动到角位置 $\theta_{1}$ 时半球速度大小
        \begin{equation}
            V _ {M} = \sqrt {\frac {1 0 m ^ {2} (R + r) g (1 - \cos \theta_ {1}) \cos^ {2} \theta_ {1}}{\left[ 7 (M + m) - 5 m \cos^ {2} \theta_ {1} \right] (M + m)}}
            \label{eq:4.p96}
        \end{equation}
        或
        \begin{equation*}
            V _ {M} ^ {2} = \frac {1 0 m ^ {2} g (R + r) \left(1 - \cos \theta_ {1}\right) \cos^ {2} \theta_ {1}}{\left[ 7 M + \left(5 \sin^ {2} \theta_ {1} + 2\right) m \right] (m + M)}
        \end{equation*}
        将上式两边对时间 $t$ 微商得
        \begin{equation*}
            2 V _ {M} a _ {M} = \frac {1 0 m g (- 2 \cos \theta + 3 \cos^ {2} \theta) [ 7 (M + m) - 5 m \cos^ {2} \theta ] - 1 0 0 m ^ {2} g (1 - \cos \theta) \cos^ {3} \theta}{[ 7 (M + m) - 5 m \cos^ {2} \theta ] ^ {2}}. \frac {m (R + r) \sin \theta \cdot \dot {\theta}}{M + m}
        \end{equation*}
        由\eqref{eq:1.p96}式可知
        \begin{equation*}
            \frac {m (R + r) \dot {\theta}}{M + m} = \frac {V _ {M}}{\cos \theta}
        \end{equation*}
        由以上两式得，小球运动到角位置 $\theta_{1}$ 时，半球的加速度大小为
        \begin{equation}
            a _ {M} \left(\theta_ {1}\right) = - \frac {5 m g \sin \theta_ {1} \left[ 1 4 (M + m) - 2 1 (M + m) \cos \theta_ {1} + 5 m \cos^ {3} \theta_ {1} \right]}{\left[ 7 (M + m) - 5 m \cos^ {2} \theta_ {1} \right] ^ {2}}
            \label{eq:5.p96}
        \end{equation}
        [（解法二）见下面第（2）问解答，列出动力学方程和运动学约束，也可以解得半球的加速度大小 $a_{M}(\theta_{1})$ 。]
    \end{subsolution}
    \begin{subsolution}
        当小球纯滚动到角位置 $\theta (\theta \leq \theta_{2})$ 时，设小球对半球的正压力和摩擦力大小分别为 $N$ 和 $f$ ，由牛顿第二定律有
        \begin{equation}
            N \sin \theta - f \cos \theta = M a _ {M}
            \label{eq:6.p96}
        \end{equation}
        在半球参考系中，对小球利用质心运动定理得
        \begin{equation}
            m g \cos \theta - N - m a _ {M} \sin \theta = m \frac {v _ {\mathrm{C}} ^ {2}}{R + r}
            \label{eq:7.p96}
        \end{equation}
        \begin{equation}
            m g \sin \theta + m a _ {M} \cos \theta - f = m \frac {\mathrm{d} v _ {\mathrm{c}}}{\mathrm{d} t}
            \label{eq:8.p96}
        \end{equation}
        式中 $v_{c}$ 为半球参考系中小球质心速度的大小
        \begin{equation}
            v _ {\mathrm{C}} = r \omega = (R + r) \dot {\theta} = \frac {M + m}{m \cos \theta} V _ {M}
            \label{eq:9.p96}
        \end{equation}
        \eqref{eq:9.p96}式的最后等式已应用了\eqref{eq:1.p96}式。在小球质心参考系中对小球利用转动定理有
        \begin{equation}
            f r = I \frac {\mathrm{d} \omega}{\mathrm{d} t}
            \label{eq:10.p96}
        \end{equation}
        由\eqref{eq:6.p96}\eqref{eq:7.p96}\eqref{eq:8.p96}\eqref{eq:9.p96}\eqref{eq:10.p96}式得
        \begin{equation}
            f = \frac {2 m}{7} \left(g \sin \theta + a _ {M} \cos \theta\right)
            \label{eq:11.p96}
        \end{equation}
        \begin{equation}
            N = m g \cos \theta - m a _ {M} \sin \theta - m \frac {v _ {\mathrm{C}} ^ {2}}{R + r}
            \label{eq:12.p96}
        \end{equation}
        按照纯滚动条件，要求
        \begin{equation*}
            f \leq \mu N
        \end{equation*}
        当小球纯滚动到角位置 $\theta_{2}$ 时开始相对于半球滑动，上式中等号成立。将\eqref{eq:11.p96}\eqref{eq:12.p96}式代入 $f = \mu N$ 得
        \begin{equation*}
            \frac {2 m}{7} \left[ g \sin \theta_ {2} + a _ {M} \left(\theta_ {2}\right) \cos \theta_ {2} \right] = \mu \left[ m g \cos \theta_ {2} - m a _ {M} \left(\theta_ {2}\right) \sin \theta_ {2} - m \frac {v _ {\mathrm{C}} ^ {2} \left(\theta_ {2}\right)}{R + r} \right]
        \end{equation*}
        将\eqref{eq:9.p96}式代入上式得 $\theta_{2}$ 所满足的方程为
        \begin{equation}
            \frac {2}{7} g \sin \theta_ {2} - \mu g \cos \theta_ {2} + a _ {M} (\theta_ {2}) \left(\frac {2}{7} \cos \theta_ {2} + \mu \sin \theta_ {2}\right) + \frac {\mu (M + m) ^ {2} V _ {M} ^ {2} (\theta_ {2})}{(R + r) m ^ {2} \cos^ {2} \theta_ {2}} = 0
            \label{eq:13.p96}
        \end{equation}
        式中 $V_{M}(\theta_{2})$ 和 $a_{M}(\theta_{2})$ 如\eqref{eq:4.p96}\eqref{eq:5.p96}式（ $\theta_1\to \theta_2$ ）所示。
    \end{subsolution}
    \begin{subsolution}
        在小球刚好运动到角位置 $\theta_{3}$ 处脱离半球的瞬间，
        \begin{equation}
            N = 0
            \label{eq:14.p96}
        \end{equation}
        此时半球的加速度为零。因此，在小球脱离半球的瞬间，小球质心相对于半球运动速度的大小 $v_{m}(\theta_{3})$ 满足
        \begin{equation}
            m g \cos \theta_ {3} = m \frac {v _ {\mathrm{C}} ^ {\prime 2}}{R + r}
            \label{eq:15.p96}
        \end{equation}
        由此得
        \begin{equation}
            v _ {\mathrm{c}} ^ {\prime} = \sqrt {(R + r) g \cos \theta_ {3}}
            \label{eq:16.p96}
        \end{equation}
    \end{subsolution}
\end{solution}